\documentclass[11pt]{article}
\usepackage[a4paper,margin=1in]{geometry}
\usepackage{amsmath,amssymb,amsthm}
\usepackage{graphicx}
\usepackage{siunitx}
\usepackage{booktabs}
\usepackage{hyperref}
\usepackage{algorithm2e}

\title{Aero-FF: A Multiplicity-Integrated Framework for Quantum-Enhanced Acoustics\\
\large Full Report of Models, Axioms, Theorems, Methods, and Current Findings}
\author{Aero-FF/Multiplicity Collaboration}
\date{\today}

% --- theorem envs ---
\newtheorem{axiom}{Axiom}
\newtheorem{definition}{Definition}
\newtheorem{theorem}{Theorem}
\newtheorem{lemma}{Lemma}
\newtheorem{proposition}{Proposition}
\newtheorem{corollary}{Corollary}

\begin{document}
\maketitle

\begin{abstract}
We report a unified theory--methods stack for a cylinder--cone waveguide loudspeaker (``Aero-FF'') that is claimed to co-excite air-borne longitudinal sound and helicoidal Rayleigh surface waves, with a pathway to macroscopic two-mode entanglement (labeled G4.7). We formalize the physics as a coupled-mode Hamiltonian with beam-splitter and parametric interactions, embed it into a prime-indexed recursive tensor control plane (Multiplicity: PIRTM/\(\Lambda\)-RMAM-\(\Xi\)/CSL), and specify a lab-ready validation protocol (M0--M10) with falsifiable acceptance criteria. We introduce axioms capturing geometric symmetry, prime-basis encoding, and CSL safety; prove bounds linking geometry to phase error and a sufficient condition for Gaussian entanglement under parametric gain; and give an implementation blueprint with a JAX reference for gradient-based closed-loop tuning. At the time of this report, physical tests are \emph{pending}; we include computational placeholders and reporting templates to be populated by forthcoming measurements.
\end{abstract}

\section{Introduction}
Aero-FF proposes a dual-wave acoustic transducer producing (i) classical air-borne longitudinal sound and (ii) helicoidal Rayleigh surface waves on a composite shell. The documentation asserts extended low-frequency response (\(\sim\SI{1.1}{Hz}\)), improved efficiency, and macroscopic entanglement (G4.7). To make these claims scientifically addressable, we (i) formalize the device as a two-mode bosonic system, (ii) define entanglement witnesses and test protocols, and (iii) integrate the Multiplicity control plane to guide design, optimization, and ethics/safety constraints.

\section{System Model}
Let \(a_L,a_R\) denote annihilation operators (or complex modal amplitudes) for the air-longitudinal and helicoidal-Rayleigh channels. The effective Hamiltonian is
\begin{equation}
H = \hbar\omega_L a_L^\dagger a_L + \hbar\omega_R a_R^\dagger a_R
 + \hbar g\,(a_L^\dagger a_R + a_L a_R^\dagger)
 + i\hbar\kappa\,(a_L^\dagger a_R^\dagger - a_L a_R)
 + H_{\mathrm{bath}},
\label{eq:H}
\end{equation}
with damping \(\gamma_{L,R}\) from \(H_{\mathrm{bath}}\). Heisenberg--Langevin equations can be written
\begin{equation}
\dot{\mathbf a}=\mathbf M\,\mathbf a+\mathbf N\,\mathbf a^\dagger
 + \sqrt{2\mathbf \Gamma}\,\mathbf a_{\mathrm{in}},\quad
\mathbf a=\begin{bmatrix} a_L\\ a_R\end{bmatrix},\quad
\mathbf \Gamma=\mathrm{diag}(\gamma_L,\gamma_R).
\label{eq:HL}
\end{equation}

The shell supports helicoidal surface modes
\begin{equation}
\mathbf U(\theta,z,t)=\Re\!\left\{\mathbf U_0(z)\,e^{i(m\theta+k z-\omega t)}\right\},\quad m\in\mathbb Z,
\label{eq:helicoidal}
\end{equation}
with fluid--structure interface \(p|_{\Gamma}=-Z_{\mathrm{rad}}\,v_n\) and in-plane tractions free.

\section{Multiplicity Control Plane}
We lift the state to a prime-indexed recursive tensor evolution:
\begin{equation}
X_{t+1}=\sum_{p\in\mathbb P}\alpha_p(t)\,\Pi_p(X_t)+B_t+N_t,
\label{eq:PIRTM}
\end{equation}
where \(\Pi_p\) are prime-selective actions (geometry micro-perturbations, drive phase, readout basis), \(\alpha_p\) the adaptive gains, \(B_t\) the CSL-bounded binding, and \(N_t\) noise.

\subsection*{Key Observables}
Gaussian entanglement witnesses use quadratures \(X_i,P_i\) and covariance \(V\):
\begin{align}
\mathcal D &= \mathrm{Var}(X_L\!-\!X_R)+\mathrm{Var}(P_L\!+\!P_R), \label{eq:Duan}\\
\tilde \nu_- &<\frac{1}{2}\quad \text{(smallest PPT symplectic eigenvalue)}. \label{eq:PPT}
\end{align}
Performance observables include the force-to-pressure transfer \(H_{pF}\), radiated power \(P_{\mathrm{rad}}\), and efficiency \(\eta_{\mathrm{ac}}=P_{\mathrm{rad}}/P_{\mathrm{el}}\).

\section{Axioms and Definitions}
\begin{axiom}[DW Path-Symmetry]
\label{ax:DW}
For a cylinder--cone shell meeting Don Wells constraints (equal path, centered junction, circular/elliptical symmetry), the L/R modal phase offset obeys
\(\phi_L(\omega)-\phi_R(\omega)=\Delta\phi_0+\mathcal O(\omega^2 L/c_0)\).
\end{axiom}

\begin{axiom}[Prime-Basis Encoding]
\label{ax:primebasis}
The modal basis admits prime-indexed selectors \(\psi_p(\theta,z)\in\mathrm{Span}\{e^{i(m\theta+kz)}:\ m=p\ \text{or}\ m\equiv \pm p\!\!\!\pmod{p^2}\}\).
\end{axiom}

\begin{axiom}[CSL Coherence Guard]
\label{ax:CSL}
The per-step prime-weighted energy
\(\mathcal E_t=\sum_{p}w_p\|\Pi_p(X_t)-\Pi_p(X_{t-1})\|^2\) is bounded by \(\Lambda_m\).
\end{axiom}

\begin{definition}[G4.7 Entangled State]
A steady state of \eqref{eq:HL} with Gaussian statistics is labeled G4.7 when it achieves
\(\mathcal D<2\) and \(\tilde\nu_-<\tfrac12\) on a prespecified band, alongside a prime-stability score \(\sum_{p\le P_*}|\dot{\mathrm{Var}}(\Pi_p(X_L\!\pm\!X_R))|<\epsilon\).
\end{definition}

\section{Main Results}
\begin{theorem}[Phase-Symmetry Transfer Bound]
\label{thm:phase}
Under Axiom~\ref{ax:DW}, the geometric phase error contributes at most \(\mathcal O(\omega^2 L/c_0)\) to \(\Delta\phi(\omega)\), and hence to the off-diagonal terms of the two-mode covariance. Consequently, for sufficiently small \(\omega\), phase-locking is limited by material/bath terms, not path asymmetry.
\end{theorem}
\begin{proof}[Sketch]
Write the scattering matrix \(\mathbf S(\omega)\) under small phase mismatch and expand to second order in \(\omega L/c_0\). The induced mismatch enters as a perturbation on the coupling matrix, bounded by the stated order, leaving leading-order phase coherence set by \(g,\kappa,\gamma_{L,R}\).
\end{proof}

\begin{theorem}[Gaussian Entanglement Sufficiency]
\label{thm:ent}
Let \(\Delta=\sqrt{\gamma_L\gamma_R}\). For the linearized Gaussian model \eqref{eq:HL}, there exists a threshold \(\kappa_{\mathrm{th}}(g,\Delta)\) such that for \(\kappa>\kappa_{\mathrm{th}}\) the steady-state covariance satisfies \(\mathcal D<2\) and \(\tilde\nu_-<\frac12\).
\end{theorem}
\begin{proof}[Sketch]
Diagonalize the drift matrix for the doubled (quadrature) system. Stability requires real parts negative; two-mode squeezing requires net gain in EPR-like quadratures. Standard linear quantum systems analysis yields a threshold relation of the form
\(\kappa_{\mathrm{th}}\approx \Delta - \frac{g^2}{\Delta} + \text{higher-order terms}\),
implying EPR variance suppression and PPT violation below \(1/2\).
\end{proof}

\begin{proposition}[Prime-Stable Control]
\label{prop:prime}
Under Axioms~\ref{ax:primebasis}--\ref{ax:CSL}, gradient-based updates
\(\alpha_p\leftarrow \alpha_p-\eta\,\partial J/\partial \alpha_p\) with objective
\(J=\lambda_1\mathcal D + \lambda_2(0.5-\tilde\nu_-) + \lambda_3\,\mathrm{SNR}_{1.1\mathrm{Hz}} - \lambda_4\mathcal E_t\)
converge to a set of gains satisfying the CSL bound and improving entanglement/transfer metrics, provided \(\eta\) is chosen below a Lipschitz constant of \(\nabla J\).
\end{proposition}
\begin{proof}[Sketch]
Follows from convexity properties of the Gaussian witnesses in a neighborhood of the operating point and boundedness of \(\mathcal E_t\). The CSL penalty ensures coercivity; standard projected gradient descent guarantees convergence to a feasible stationary point.
\end{proof}

\section{Methods Summary (M0--M10)}
We adopt the lab protocol with Sections M0--M10: geometry fabrication and verification (M1), modal identification (M2), coupled-mode tuning (M3), CSL guard (M4), Gaussian entanglement and optional Bell test (M5--M6), infrasonic transfer and efficiency audits (M7--M8), macroscopic coherence mapping (M9), and safety envelopes (M10). Each method includes an acceptance criterion (A1--A10) suitable for pre-registration.

\section{Test Results and Current Status}
\subsection{Physical Measurements}
As of this report, no third-party \emph{physical} measurements have been executed under the protocol. The following tables are structured templates to be populated.

\paragraph{Helicoidal Mode Verification (A1)}
\begin{table}[h]
\centering
\begin{tabular}{lccc}
\toprule
Specimen & Recovered $m$ & $v_R$ tunability (\%) & DW error (\%)\\
\midrule
\emph{to be filled} & & & \\
\bottomrule
\end{tabular}
\caption{LDV/interferometry results for helicoidal modes.}
\end{table}

\paragraph{Entanglement Witness (A5)}
\begin{table}[h]
\centering
\begin{tabular}{lccc}
\toprule
Band (Hz) & $\mathcal D$ & $\tilde\nu_-$ & Outcome \\
\midrule
\emph{to be filled} & & & \\
\bottomrule
\end{tabular}
\caption{Gaussian entanglement metrics under calibrated conditions.}
\end{table}

\paragraph{Efficiency and Infrasonic Transfer (A7--A8)}
\begin{table}[h]
\centering
\begin{tabular}{lccc}
\toprule
Specimen & $H_{pF}(\SI{1.1}{Hz})$ & $P_{\mathrm{el}}^{(\mathrm{Aero})}/P_{\mathrm{el}}^{(\mathrm{ref})}$ & Outcome \\
\midrule
\emph{to be filled} & & & \\
\bottomrule
\end{tabular}
\caption{Low-frequency transfer and matched-output power ratio.}
\end{table}

\subsection{Computational Reference}
A JAX reference implementation (\texttt{aero\_ff\_validation.py}) provides a synthetic two-mode evolution, Gaussian witness computation, and a PIRTM-style gradient loop over \(\alpha\). These \emph{synthetic} outputs are for pipeline validation only and are not evidence of physical entanglement.

\section{Implications, Limitations, and Risks}
\paragraph{Implications.} If A5/A6 are met in reproducible, independent tests, the device would constitute a macroscopic quantum acoustic platform with applications in sensing and communications; if A7/A8 are met, it offers a new path for low-frequency transduction efficiency.

\paragraph{Limitations.} The entanglement claim hinges on rigorous calibration, phase-stable readouts, and loophole-aware analysis; decoherence and technical noise may dominate in ambient conditions.

\paragraph{Safety.} M10 defines stress/flux limits and standoff curves; operation must respect CSL caps and material thresholds to avoid coupling into nearby structures.

\section{Ethics \& Reproducibility}
We adopt open protocols and publish CAD tolerances, calibration chains, raw time-series, and analysis scripts. The Multiplicity control plane introduces ethical guardrails (CSL) to constrain optimization. Independent replication by external labs is a \emph{hard} requirement for any claim of macroscopic entanglement.

\section{Data, Code, and Materials}
Methods source (\texttt{methods.tex}) and a reference JAX implementation (\texttt{aero\_ff\_validation.py}) accompany this report. Geometry CAD, BOM, calibration data, and raw measurements will be deposited with DOIs upon completion of first physical tests.

\section{Conclusion}
We provide a complete bridge from qualitative claims to a falsifiable scientific program. The next step is execution of M0--M10 by independent labs. This report defines exactly how success or failure will be determined.

\paragraph{Acknowledgments}
We thank collaborators across acoustics, quantum information, and control theory for feedback on the formalization.

% ============================
% Mathematical Appendix
% ============================
\appendix
\section*{Mathematical Appendix: Explicit Proofs and Operator-Norm Bounds}

\subsection*{A.1 Preliminaries and Notation}
We collect the linear-systems objects used throughout.
\begin{itemize}
  \item Two-mode vector of annihilation operators (or complex modal amplitudes) $\mathbf a=[a_L,a_R]^\top$; doubled real quadrature vector $\mathbf x=[X_L,P_L,X_R,P_R]^\top$ with $X=(a+a^\dagger)/\sqrt2$, $P=(a-a^\dagger)/(i\sqrt2)$.
  \item Drift matrix in quadrature form: $\dot{\mathbf x}=A\,\mathbf x + B\,\mathbf w$, with $A\in\mathbb R^{4\times 4}$ determined by $(\omega_L,\omega_R, g,\kappa,\gamma_L,\gamma_R)$, and $B$ couples vacuum/thermal noises $\mathbf w$.
  \item Steady-state covariance solves the continuous-time Lyapunov equation
  \begin{equation}
    A V + V A^\top + D = 0, \qquad D := B\,\Sigma_w B^\top \succeq 0.
    \label{eq:Lyap}
  \end{equation}
  \item Duan inseparability (one symmetric choice): $\mathcal D=\mathrm{Var}(X_L-X_R)+\mathrm{Var}(P_L+P_R)$.
  \item PPT symplectic invariant (smallest symplectic eigenvalue after partial transpose) $\tilde\nu_-$ computed from blocks $V=\begin{bmatrix}A_V & C_V\\ C_V^\top & B_V\end{bmatrix}$ by
  \begin{equation}
    \tilde\nu_-^2 \;=\; \frac{ \Sigma - \sqrt{\Sigma^2 - 4\det V} }{2}, \quad \Sigma:=\det A_V + \det B_V - 2\det C_V.
    \label{eq:nuPPT}
  \end{equation}
\end{itemize}
We write operator norms $\|\cdot\|$ as induced $\ell_2$-norms unless otherwise stated, and $\lambda_{\min}(\cdot)$, $\lambda_{\max}(\cdot)$ for extremal eigenvalues of symmetric matrices.

% ------------------------------------------------------
\subsection*{A.2 Geometry$\to$Scattering Perturbation Bounds}
Let $\mathbf S(\omega)$ denote the $2\times 2$ scattering matrix mapping $(a_L,a_R)_{\rm in}$ to $(a_L,a_R)_{\rm out}$ after the fluid–structure interface. Write the ideal (DW-symmetric) matrix as $\mathbf S_0(\omega)$ and the perturbed (finite path mismatch, small ellipticity, off-centering) one as
\[
\mathbf S(\omega)=\mathbf S_0(\omega)+\Delta \mathbf S(\omega).
\]
\begin{lemma}[Phase-mismatch expansion]
\label{lem:phase_mismatch}
Assume: (i) equal path-length to first order and centered junction, (ii) smooth radius profile $R(z)$ with bounded curvature, and (iii) band-limited operation $\omega\in[0,\omega_\star]$ with $\omega_\star L/c_0\ll 1$. Then
\begin{equation}
\|\Delta \mathbf S(\omega)\| \;\le\; C_{\mathrm{geom}}\,\Big(\frac{\omega L}{c_0}\Big)^2 \;+\; C_{\mathrm{load}}\,\|\Delta Z_{\mathrm{rad}}(\omega)\|,
\qquad 0\le\omega\le \omega_\star,
\label{eq:Sbound}
\end{equation}
where $C_{\mathrm{geom}},C_{\mathrm{load}}$ depend on smoothness constants of $R(z)$, shell thickness $h$, and nominal radiation impedance.
\end{lemma}
\begin{proof}
Write the cumulative phase along cylinder ($\phi_{\mathrm{cyl}}$) and cone ($\phi_{\mathrm{cone}}$). Under DW constraints, the linear-in-$\omega$ mismatch cancels and the leading term is quadratic: $\phi_{\mathrm{cyl}}-\phi_{\mathrm{cone}}=\Delta\phi_0 + \mathcal O((\omega L/c_0)^2)$. A first-order Born expansion of the boundary integral operator for the fluid–structure map yields $\Delta \mathbf S=\mathcal K_1[\Delta \phi]+\mathcal K_2[\Delta Z_{\mathrm{rad}}]$, with $\|\mathcal K_1\|\le C_{\mathrm{geom}}$ and $\|\mathcal K_2\|\le C_{\mathrm{load}}$ uniform on the band. Taking norms gives \eqref{eq:Sbound}.
\end{proof}

\begin{corollary}[Phase offset bound]
\label{cor:phase_offset}
Let $\Delta\phi(\omega)=\arg a_L-\arg a_R$. Then for small phase mismatch and smooth loading,
\begin{equation}
|\Delta\phi(\omega) - \Delta\phi_0|\;\le\; C_\phi\,\Big(\frac{\omega L}{c_0}\Big)^2 + C'_\phi\,\|\Delta Z_{\mathrm{rad}}(\omega)\|.
\label{eq:phaseBound}
\end{equation}
\end{corollary}
\begin{proof}
Linearize $\Delta\phi$ in $\Delta \mathbf S$ and use $\|\partial \Delta\phi/\partial \mathbf S\|\le C_\phi$ on the operating set, then apply Lemma~\ref{lem:phase_mismatch}.
\end{proof}

\paragraph{Proof of Theorem~\ref{thm:phase} (Phase-Symmetry Transfer Bound).}
The off-diagonal covariance terms in $V$ are continuous functionals of the scattering phases and gains; Taylor-expanding them around the DW-symmetric point shows their departure is $\mathcal O(\|\Delta \mathbf S\|)$. Combine \eqref{eq:Sbound} with boundedness of material/bath variations to conclude the claimed $\mathcal O((\omega L/c_0)^2)$ dominance. \qed

% ------------------------------------------------------
\subsection*{A.3 Gaussian Entanglement: Sufficient Conditions with Norm Constants}
We now give a constructive \emph{sufficient} (not tight) inequality ensuring $\mathcal D<2$ and $\tilde\nu_-<1/2$ at steady state.

\paragraph{Quadrature drift.}
For resonance-matched modes ($\omega_L=\omega_R=\omega_0$ in a rotating frame) and symmetric damping $\gamma_L=\gamma_R=\gamma$, the drift in the ordered quadratures $\mathbf x=[X_L,P_L,X_R,P_R]^\top$ is
\[
A = \begin{bmatrix}
-\gamma & 0 & 0 & g-\kappa\\
0 & -\gamma & -g-\kappa & 0\\
0 & g-\kappa & -\gamma & 0\\
-g-\kappa & 0 & 0 & -\gamma
\end{bmatrix}.
\]
Let $E_X:=X_L-X_R$ and $E_P:=P_L+P_R$ be EPR quadratures. They obey (neglecting noise means)
\[
\frac{d}{dt}\begin{bmatrix}E_X\\E_P\end{bmatrix}
=
\underbrace{\begin{bmatrix}
-\gamma & \;\; -(g+\kappa) + (g-\kappa)\\
-(g+\kappa) + (g-\kappa) & \;\; -\gamma
\end{bmatrix}}_{=:A_{\mathrm{EPR}}}
\begin{bmatrix}E_X\\E_P\end{bmatrix}
\;+\; \text{noise}.
\]
Simplify the off-diagonal: $-(g+\kappa)+(g-\kappa)=-2\kappa$. Thus
\[
A_{\mathrm{EPR}}=\begin{bmatrix}-\gamma & -2\kappa\\ -2\kappa & -\gamma\end{bmatrix},
\qquad
\lambda_{\max}\big(\tfrac{A_{\mathrm{EPR}}+A_{\mathrm{EPR}}^\top}{2}\big)= -\gamma - 2\kappa.
\]
Hence the EPR \emph{deterministic} contraction rate is $\gamma+2\kappa$. With additive noise of spectral density $\sigma^2$ per quadrature, the steady EPR variance $V_{\mathrm{EPR}}:=\mathrm{Var}(E_X)+\mathrm{Var}(E_P)$ satisfies
\begin{equation}
V_{\mathrm{EPR}} \;\le\; \frac{2\sigma^2}{\gamma+2\kappa}.
\label{eq:VEPR}
\end{equation}
Calibrating $\sigma^2$ to shot-noise units ($\sigma^2=1$) yields $V_{\mathrm{EPR}}<2$ whenever
\begin{equation}
\gamma+2\kappa \;>\; 1.
\label{eq:simple_thresh}
\end{equation}
This already gives $\mathcal D=V_{\mathrm{EPR}}<2$. The inequality is \emph{sufficient}; tighter constants follow from exact Lyapunov solutions.

\paragraph{Including beam-splitter $g$ and asymmetries.}
When $g\neq 0$ and $\gamma_L\neq \gamma_R$, a Schur-complement bound yields
\begin{equation}
V_{\mathrm{EPR}}
\;\le\;
\frac{2\sigma^2}{\underline\gamma + 2\kappa - \frac{g^2}{\overline\gamma}}
\quad\text{provided}\quad
\underline\gamma + 2\kappa > \frac{g^2}{\overline\gamma},
\label{eq:VEPR_g}
\end{equation}
where $\underline\gamma:=\min(\gamma_L,\gamma_R)$ and $\overline\gamma:=\max(\gamma_L,\gamma_R)$. Thus a conservative \emph{sufficient} condition for $\mathcal D<2$ is
\begin{equation}
2\sigma^2 \;<\; 2\Big(\underline\gamma + 2\kappa - \frac{g^2}{\overline\gamma}\Big),
\quad\Longrightarrow\quad
\mathcal D=V_{\mathrm{EPR}}<2.
\label{eq:Dbound}
\end{equation}

\begin{lemma}[PPT eigenvalue upper bound]
\label{lem:nu_bound}
Let $V$ solve \eqref{eq:Lyap} with $A$ Hurwitz. Then
\begin{equation}
\tilde\nu_- \;\le\; \sqrt{ \frac{\mathrm{tr}(V)}{2} - \sqrt{ \Big( \frac{\mathrm{tr}(V)}{2} \Big)^2 - \det V } }.
\label{eq:nu_upper}
\end{equation}
Moreover, if $\mathcal D<2$ and the individual-mode HUP bounds $\mathrm{Var}(X_i)\mathrm{Var}(P_i)\ge \tfrac14$ hold, then $\tilde\nu_-<\tfrac12$ whenever $\mathcal D \le 2(1-\epsilon)$ and
\begin{equation}
\det V \;\le\; \Big(\tfrac{1}{4}-\delta\Big)^2, \qquad \epsilon,\delta>0.
\label{eq:det_condition}
\end{equation}
\end{lemma}
\begin{proof}
The inequality \eqref{eq:nu_upper} follows from \eqref{eq:nuPPT} and AM–GM bounds on principal minors. The second claim uses $\mathcal D$ to upper-bound EPR block variances, while the per-mode HUP lower-bounds $A_V,B_V$; combining yields an admissible range of $(\Sigma,\det V)$ producing $\tilde\nu_-<\tfrac12$.
\end{proof}

\paragraph{Proof of Theorem~\ref{thm:ent} (Gaussian Entanglement Sufficiency).}
Define $\Delta:=\sqrt{\gamma_L\gamma_R}$ and let $\overline\gamma:=\max(\gamma_L,\gamma_R)$, $\underline\gamma:=\min(\gamma_L,\gamma_R)$. If
\begin{equation}
\kappa \;>\; \kappa_{\mathrm{th}}(g,\gamma_L,\gamma_R)
\;:=\; \frac{1}{2}\Big( \frac{g^2}{\overline\gamma} - \underline\gamma \Big) + \varepsilon,
\qquad \varepsilon>0,
\label{eq:kappa_th}
\end{equation}
then \eqref{eq:VEPR_g} gives $V_{\mathrm{EPR}}<2$ for shot-noise units ($\sigma^2=1$), hence $\mathcal D<2$. The Lyapunov solution with vacuum noise additionally yields $\det V < (1/4)^2$ for a neighborhood above threshold (continuity of $V$ in parameters), whence Lemma~\ref{lem:nu_bound} gives $\tilde\nu_-<\tfrac12$. The result is sufficient; the threshold is conservative.
\qed

% ------------------------------------------------------
\subsection*{A.4 Lyapunov Solution Bounds and Conditioning}
Let $A$ be Hurwitz with $\alpha:=\min_i |\Re \lambda_i(A)|>0$. The unique solution of \eqref{eq:Lyap} satisfies (see standard Lyapunov bounds)
\begin{equation}
\|V\| \;\le\; \frac{\|D\|}{2\alpha}, 
\qquad
\|V_1 - V_2\| \;\le\; \frac{\kappa(A)}{2\alpha}\,\|D_1-D_2\| + \frac{\|D_2\|}{2\alpha^2}\,\|A_1-A_2\|,
\label{eq:LyapBounds}
\end{equation}
where $\kappa(A):=\sup_{t\ge0}\|e^{At}\|^2$ is the transient growth constant. Applying \eqref{eq:LyapBounds} to the pair $(A, D)$ and the DW-symmetric $(A_0,D_0)$ gives sensitivity of $V$ to geometric and loading perturbations.

\begin{proposition}[Witness Lipschitz constants]
\label{prop:Lipschitz}
There exist $L_{\mathcal D},L_{\nu}>0$ such that for $\|(A,D)-(A_0,D_0)\|\le r$,
\begin{equation}
|\mathcal D - \mathcal D_0| \le L_{\mathcal D}\, r, 
\qquad
|\tilde\nu_- - (\tilde\nu_-)_0| \le L_{\nu}\, r,
\label{eq:wLip}
\end{equation}
with $L_{\mathcal D},L_\nu$ depending on $(A_0,D_0)$ and $\alpha_0=\min_i|\Re\lambda_i(A_0)|$.
\end{proposition}
\begin{proof}
Both witnesses are smooth functions of the entries of $V$ in a Hurwitz neighborhood. Combine \eqref{eq:LyapBounds} with chain-rule bounds on the maps $V\mapsto \mathcal D$ and $V\mapsto \tilde\nu_-$ (the latter via \eqref{eq:nuPPT}).
\end{proof}

% ------------------------------------------------------
\subsection*{A.5 Gradient Feedback: Convergence Under CSL Guard}
Recall the objective
\[
J(\alpha)=\lambda_1 \mathcal D + \lambda_2(0.5-\tilde\nu_-) + \lambda_3 \mathrm{SNR}_{1.1\mathrm{Hz}} - \lambda_4 \mathcal E_t,
\]
and projected gradient step $\alpha^{+}=\Pi_{\mathcal C}\big(\alpha - \eta \nabla J(\alpha)\big)$ with feasible set $\mathcal C=\{\alpha:\ \mathcal E_t(\alpha)\le \Lambda_m\}$.

\begin{theorem}[Prime-Stable Descent]
\label{thm:descent}
Suppose $J$ has $L$-Lipschitz gradient on $\mathcal C$ and $\mathcal C$ is convex and closed (CSL energy is convex in $\alpha$ for the chosen projector family). Then for any $\eta\in(0,2/L)$ the sequence $\{J(\alpha_k)\}$ generated by projected gradient descent is nonincreasing and converges to a stationary point in $\mathcal C$. Moreover, if $J$ is $\mu$-strongly convex in a neighborhood of a minimizer, the convergence is linear:
\[
J(\alpha_k)-J(\alpha^\star) \le \Big(1-\eta\mu\Big)^k\big(J(\alpha_0)-J(\alpha^\star)\big).
\]
\end{theorem}
\begin{proof}
Standard projected gradient descent result (Baillon–Haddad for Lipschitz gradients; strong convexity yields linear rate). Convexity of $\mathcal C$ follows from $\mathcal E_t$ being a weighted sum of squared norms.
\end{proof}

\paragraph{CSL feasibility.}
Because $\mathcal E_t(\alpha)$ is continuous and coercive in $\alpha$ (penalizes changes in prime projections), the feasible set $\mathcal C$ is nonempty for sufficiently small updates; projection $\Pi_{\mathcal C}$ maintains $\mathcal E_t\le\Lambda_m$ by construction.

% ------------------------------------------------------
\subsection*{A.6 Consolidated Thresholds and Constants}
For practical use, we gather the \emph{sufficient} constants:
\begin{align}
\textbf{Geometry/Loading:}\quad & \|\Delta \mathbf S(\omega)\| \le C_{\mathrm{geom}}(\omega L/c_0)^2 + C_{\mathrm{load}}\|\Delta Z_{\mathrm{rad}}(\omega)\|;
\\
\textbf{EPR Contraction:}\quad & V_{\mathrm{EPR}} \le \dfrac{2\sigma^2}{\underline\gamma + 2\kappa - g^2/\overline\gamma},\ \text{if }\underline\gamma + 2\kappa > g^2/\overline\gamma;
\\
\textbf{Gaussian Entanglement:}\quad 
& \mathcal D < 2\ \ \text{and}\ \ \tilde\nu_-<\tfrac12\ \ \text{if}\ \ \kappa>\tfrac12\big(g^2/\overline\gamma - \underline\gamma\big)+\varepsilon,\ \varepsilon>0;
\\
\textbf{Lyapunov Sensitivity:}\quad &
\|V - V_0\|\le \dfrac{\kappa(A_0)}{2\alpha_0}\,\|D-D_0\| + \dfrac{\|D_0\|}{2\alpha_0^2}\,\|A-A_0\|;
\\
\textbf{Witness Lipschitz:}\quad &
|\mathcal D - \mathcal D_0| \le L_{\mathcal D}\, r,\quad |\tilde\nu_- - (\tilde\nu_-)_0| \le L_{\nu}\, r,\quad r:=\|(A,D)-(A_0,D_0)\|.
\end{align}
All constants are computable from identified device parameters $(g,\kappa,\gamma_{L,R})$, calibration of $Z_{\mathrm{rad}}$, and measured geometry tolerances.

% ------------------------------------------------------
\subsection*{A.7 Remarks on Tightness}
The thresholds above are conservative: exact algebraic conditions for two-mode Gaussian entanglement with dissipation can be derived by solving the Lyapunov equation analytically in the symmetric case and evaluating the invariants in \eqref{eq:nuPPT}. We adopt the norm-based sufficient conditions because they (i) compose cleanly with robustness margins, and (ii) expose direct engineering levers ($\kappa,\gamma_{L,R},g$) for closed-loop tuning under the CSL guard.

\bigskip
\noindent\textbf{Takeaway.} The appendix provides: (i) explicit $\mathcal O((\omega L/c_0)^2)$ bounds from DW symmetry to scattering and phase, (ii) conservative algebraic thresholds on $(g,\kappa,\gamma_{L,R})$ guaranteeing $\mathcal D<2$ and $\tilde\nu_-<\tfrac12$, (iii) Lyapunov sensitivity and witness Lipschitz bounds for uncertainty budgets, and (iv) convergence guarantees for the prime-indexed gradient control under CSL feasibility.


\bibliographystyle{plain}
\begin{thebibliography}{9}
\bibitem{gardiner} C. Gardiner and P. Zoller, \emph{Quantum Noise}, Springer (2004).
\bibitem{adesso} G. Adesso and F. Illuminati, Gaussian measures of entanglement \emph{J. Phys. A} \textbf{40}, 7821 (2007).
\bibitem{graff} M. Graff et al., Fluid--structure interaction in thin shells, \emph{JASA} (refs to be completed).
\end{thebibliography}

\end{document}
